\documentclass[11pt]{article}

\usepackage[preprint]{acl}
\usepackage{times}
\usepackage{latexsym}
\usepackage[T1]{fontenc}
\usepackage[utf8]{inputenc}
\usepackage{microtype}
\usepackage{graphicx}
\usepackage{booktabs}

\title{Revisiting Multi-Agent Social Reasoning for Conversational Power:\\Factor Graphs, Decision Models, and Context}

\author{Muttaqi Ahmad Alladin \\
  Purdue University \And
  Krish Patel \\
  Purdue University \And
  Jiwon Kim \\
  Purdue University}

\begin{document}
\maketitle

\begin{abstract}
\citet{wu2025power} predict which of two speakers in a six-turn meeting excerpt is the professor, with a factor graph over LLM-written reasons and rebuttals.
We proposed to improve the graph with quote-checked arguments and Dung's grounded extension.
On the 151 released snippets neither change helps, and no method we tried scores above the majority label.
Tracing every snippet back to the ICSI Meeting Corpus shows that the released text and labels disagree about who is speaking: under the reading of the speaker letters that makes the labels correct, 57.6\% of the words are displayed under the other speaker's letter.
We rebuilt the snippets from the corpus, taking speaker turns from the corpus itself.
On the 150 rebuilt snippets a System One decision model (Jev) identifies the professor in 71.3\% of cases against a 56.7\% majority baseline, and in 89.3\% when it also sees 60 segments of meeting on each side.
Their factor graph, rerun call by call with arguments from GPT-6.1 Terra and scores from Jev, reaches 72.0\%, level with asking Terra directly (72.7\%).
A factor graph built from typed yes/no statements, with no generated arguments, matches the direct question at every context length.
\end{abstract}

\section{Introduction}

\citet{wu2025power} study professor-student exchanges of six turns from the ICSI Meeting Corpus \citep{janin2003icsi} and take the professor to hold more power.
They report 46.3\% for six PhD judges, 49.0\% for GPT-3.5 on the snippet alone, and 54.3\% for a factor graph over LLM-written reasons and rebuttals solved with AD3 \citep{martins2011augmented}.
In the label file they released, speaker A is the higher-power speaker in 83 of 151 rows, so always answering A scores 55.0\%.

Our proposal kept their cached arguments and changed three things: a quote check on every argument, a rebuttal prompt that shows the rebuttal, and Dung's grounded extension \citep{dung1995acceptability} in place of their XOR program.
Section~\ref{sec:released} reports those experiments on the released snippets.
Section~\ref{sec:data} traces the released snippets back to the corpus, shows that their displayed speakers and their labels disagree, and describes a snippet set rebuilt from the corpus.
Section~\ref{sec:rebuilt} reports the same task on snippets rebuilt from the corpus, with the original factor graph, a decision model, zero-shot LLMs, and longer context.

\section{The released snippets}
\label{sec:released}

\begin{table}[t]
\centering\small
\setlength{\tabcolsep}{3pt}
\begin{tabular}{lrr}
\toprule
Method & Correct & Acc. \\
\midrule
Always A & 83 & 55.0 \\
XOR graph, cached GPT-3.5 scores & 82 & 54.3 \\
Grounded extension, as proposed & 69 & 45.7 \\
XOR graph with quote filter & 64 & 42.4 \\
Gemma 4 zero-shot & 66 & 43.7 \\
Jev, ``is X the professor'' & 72 & 47.7 \\
XOR graph, Jev scores & 76 & 50.3 \\
Preference grounded, cached scores & 79 & 52.3 \\
Preference grounded, Jev scores & 71 & 47.0 \\
Typed-statement graph (Jev) & 77 & 51.0 \\
\bottomrule
\end{tabular}
\caption{Accuracy on the 151 released snippets. Beating the 55.0\% base rate at $p<0.05$ needs 94 correct.}
\label{tab:released}
\end{table}

Table~\ref{tab:released} lists the runs.
Our reimplementation of the XOR graph as an integer program reproduces the published 54.3\% exactly (82 of 151).
The quote filter kept 33.8\% of the 7,248 arguments and lowered accuracy.

\paragraph{The grounded extension is degenerate as we specified it.}
Every reason has one paired rebuttal and nothing attacks a rebuttal, so every rebuttal is \textsc{in} and all 3,624 reasons are \textsc{out}.
All 151 decisions fall to the tie-break.
Removing weak rebuttals does not help, because we also made each reason for A and each reason for B attack each other, and under grounded semantics a mutual attack leaves both arguments undecided.
A preference-based framework, where an attack succeeds only when the attacker scores at least as high as its target, removes both problems and gives a decision from \textsc{in} counts on every snippet.
It is well behaved and no more accurate (79 of 151 with cached scores).

\paragraph{Rescoring does not help.}
A logistic regression fit to the cached validity and grounding scores, evaluated leave-one-out, reaches 53.0 to 57.0\%.
Rescoring all arguments with Jev, including the corrected rebuttal question on all 3,624 rebuttals, gives 50.3\% under the XOR graph.

\section{Data}
\label{sec:data}

\paragraph{Sources.}
The released test set is the file \texttt{meeting.csv} of \citet{wu2025power}: 151 rows, each with a conversation between speakers A and B and a label naming the higher-power speaker.\footnote{\url{https://github.com/kandicew/social-reasoning}, commit \texttt{fed11b6}, \texttt{data/meeting.csv}.}
The snippets come from the ICSI Meeting Corpus \citep{janin2003icsi}.
We use its original MRT transcripts, in which every segment carries a start time, a participant ID and the transcribed text, and its participant metadata, which records each participant's education level.\footnote{\url{https://groups.inf.ed.ac.uk/ami/icsi/download}: \texttt{ICSI\_original\_transcripts.zip} and \texttt{ICSI\_core\_NXT.zip} (\texttt{speakers.xml}), distributed under CC BY 4.0.}

\paragraph{Locating the released snippets.}
After removing transcript markup, the non-empty lines of a released snippet are a run of consecutive corpus segments.
Matching on that sequence finds all 151 snippets, each in exactly one place, across 45 meetings.
This step uses only the released text.
It gives, for every line of every snippet, the participant who said it.

\begin{table}[t]
\centering\small
\begin{tabular}{lll}
\toprule
Shown & Said by & Text \\
\midrule
A & professor & Uh, maybe you have to\ldots \\
A & other & Mm-hmm. \\
B & professor & but - \\
A & other & Yeah, I guess. \\
B & professor & Yeah. \\
A & professor & do you have any way of\ldots \\
A & other & Mmm. \\
B & other & No, we don't. \\
B & professor & Yeah. \\
\bottomrule
\end{tabular}
\caption{Released snippet 0 (meeting Bro025, starting at 1273.5\,s, abridged). ``Shown'' is the letter the line appears under in the released text. ``Said by'' is from the corpus participant IDs. The released label is B.}
\label{tab:example}
\end{table}

\paragraph{Displayed speakers against corpus speakers.}
In the released text each displayed turn ends with the first segment of the next corpus speaker, as in Table~\ref{tab:example}: the turn shown under A contains the professor's remarks and then the other participant's ``Mm-hmm.''
A letter therefore does not stand for one participant throughout a snippet, and there are two ways to read it.

If each letter is read as the speaker of the last line of its turns, the letters map one-to-one onto the two participants in 128 snippets with exactly one professor, and the released label names the professor in 128 of those 128.
Under that reading, 7,596 of the 13,185 words in those snippets (57.6\%) are displayed under the other participant's letter, and in 73 of the 128 snippets more than half of the words are.

If instead each letter is read as whoever says most of the words shown under it, the text is mostly attributed correctly, and the released label names the letter that carries most of the professor's words in only 69 of the 128 snippets.

Under either reading the text and the labels disagree: the letter that a reader sees doing the professor's talking is, in about half of the snippets, not the letter the label rewards.
The remaining 23 snippets have a third speaker, no single professor, or a letter that maps to two participants.

\paragraph{What we can and cannot conclude.}
The cached GPT-3.5 arguments in the release were generated from this text: the conversation stored with each cached snippet is identical to the corresponding row of \texttt{meeting.csv} for all 151.
The code that built the snippets from the corpus is not in the public repository, so we cannot say which step produced the mismatch.
We do not know whether the human judges of \citet{wu2025power} saw the same text, or whether the numbers in their paper were computed from this file.
We have not yet confirmed these observations with the authors.

\paragraph{Consequence for the task.}
A reader or model that judges the displayed text will, in many snippets, credit the professor's behaviour to the other letter.
This is consistent with Section~\ref{sec:released}, where every method is at or slightly below the majority rate, and with the near-chance human and LLM accuracies reported by \citet{wu2025power}.
It also predicts that the same judgments should become accurate once speech is attributed to the people who said it, which Section~\ref{sec:rebuilt} tests.

\paragraph{The rebuilt set.}
We rebuilt each snippet from the located corpus segments.
Speakers are the corpus participant IDs, consecutive segments by one participant form a turn, and the professor is the participant whose education level in \texttt{speakers.xml} is Professor.
A snippet is kept if it has exactly one professor and at least two speakers, which drops one snippet and leaves 150.
In the 128 snippets above, the letters follow the last-line reading, so the released labels carry over unchanged; in the other 22, the first of the two speakers to talk is A.
In total 142 of the 150 labels equal the released label, and the professor is A in 85 (a 56.7\% base rate).
Three snippets include a third participant, shown as C.
Each snippet is rendered at three lengths: its own segments, and 15 or 60 further corpus segments on each side, with other participants shown as C, D, and so on.
The other speaker's education level is PhD in 117 snippets, Grad in 30 and Postdoc in 3, so ``student'' follows the usage of \citet{wu2025power}.
The builder is deterministic and regenerates the set byte for byte from the three source files.

\section{Snippets rebuilt from the corpus}
\label{sec:rebuilt}

All experiments in this section use the rebuilt set of Section~\ref{sec:data}: 150 snippets with speakers taken from the corpus, at the same span as the released snippet and with 15 or 60 extra segments on each side.

\paragraph{Jev.}
Jev is a decision model that returns a probability for each typed yes/no question about a supplied state.\footnote{TypeSafe AI, \texttt{jev-1.14.0}, \url{https://docs.typesafe.ai}. The questions used on the rebuilt snippets were fixed before that data was built.}
The direct question is ``Is speaker X the professor?'', asked for A and for B, and we predict the speaker with the higher probability.

\paragraph{The factor graph of Wu et al.}
We ran their method on the rebuilt snippets with arguments from current models, in two ways.
In the first, GPT-6.1 Terra wrote the arguments call by call with their prompt templates: one independent request for each aspect summary, each list of three reasons, and each rebuttal (36 requests per snippet), with tools disabled and no access to labels.
In the second, Claude Haiku wrote all arguments for a snippet in one pass.
Jev scored both sets: validity and grounding of each argument, whether each rebuttal counters its reason, and similarity and contradiction of reason pairs.
A separate Haiku instance also scored the Haiku arguments.
Their rule file and weights (reason 20, rebuttal 18) do the inference, in the three variants of the paper.

\paragraph{Grounded extension.}
We also ran the aggregator from our proposal on the same arguments.
The framework has 48 arguments per snippet, 24 reasons and their 24 rebuttals.
A reason has strength $P(\textrm{valid})\cdot P(\textrm{grounded})$ and a rebuttal has strength $P(\textrm{counters its reason})\cdot P(\textrm{grounded})$.
In the \emph{proposed} relation every rebuttal attacks its reason, and every reason for A and every reason for B attack each other.
In the \emph{preference} relation an attack succeeds only if the attacker is at least as strong as its target: a rebuttal defeats its reason when it is no weaker, and a reason defeats an opposing reason when it is strictly stronger.
The grounded labelling marks an argument \textsc{in} when all its attackers are \textsc{out} and \textsc{out} when some attacker is \textsc{in}.
The speaker with more \textsc{in} reasons is predicted, with ties going to the larger sum of reason strengths and then to A.

\paragraph{Typed-statement factor graph.}
We also keep their graph and replace the generated text.
For each of their four aspects we wrote three yes/no statements, such as ``Does speaker X evaluate speaker Y's work?''.
Jev's probability for a statement about X is the weight of a reason for X, in the manner of \citet{kassner2023rationality}, and the same statement about the other speaker is its rebuttal.
One call with 24 questions replaces the generation and scoring calls for 48 arguments per snippet.
Unlike the argumentation scoring studied by \citet{sanayei2025debates}, no model is asked to reason over the graph.

\begin{table}[t]
\centering\small
\setlength{\tabcolsep}{3pt}
\begin{tabular}{lrrrr}
\toprule
 & Rel. & Rebuilt & $\pm$15 & $\pm$60 \\
\midrule
More words & 46.0 & 66.7 & 66.0 & 68.7 \\
Longer turns & & 68.0 & 69.3 & 66.7 \\
Jev, professor & 48.0 & 71.3 & 84.7 & 89.3 \\
Jev, more power & 48.0 & 70.7 & 79.3 & 88.0 \\
Typed-statement graph & 51.3 & 70.0 & 80.7 & 90.7 \\
\bottomrule
\end{tabular}
\caption{Accuracy (\%) on the 150 snippets. ``Rel.''\ uses the released text and labels for the same snippets. The other columns use the rebuilt text at the same span and with 15 or 60 extra segments per side. Always answering A scores 56.7\%. The blank cell was not run.}
\label{tab:rebuilt}
\end{table}

\begin{table}[t]
\centering\small
\setlength{\tabcolsep}{3pt}
\begin{tabular}{llrrrrr}
\toprule
Args. & Scores & Dir. & Reas. & +Reb. & All & Gr. \\
\midrule
Terra & Jev & 72.7 & 66.7 & 72.0 & 72.0 & 70.0 \\
Haiku & Jev & 68.0 & 73.3 & 70.0 & 70.0 & 73.3 \\
Haiku & Haiku & 68.0 & 68.7 & 67.3 & 68.0 & 64.7 \\
\bottomrule
\end{tabular}
\caption{Argument-based methods on the rebuilt snippets at the same span (accuracy, \%). Rows give who wrote and who scored the arguments. ``Dir.''\ is the argument-writing model asked directly. ``Reas.'', ``+Reb.''\ and ``All'' are the three variants of the factor graph of \citet{wu2025power}. ``Gr.''\ is the grounded extension with the preference relation. For comparison, the direct Jev question scores 71.3\% and Claude Sonnet asked directly scores 76.0\%.}
\label{tab:graph}
\end{table}

\paragraph{Results.}
Table~\ref{tab:rebuilt} gives accuracy by input for the methods that need no generated arguments.
On the rebuilt snippets at the same span, the direct Jev question identifies the professor in 107 of 150 cases.
Against always answering A this is 34 wins and 12 losses (exact McNemar $p=0.0016$), and a bootstrap that resamples whole meetings gives a 95\% interval of 6.7 to 22.0 points for the difference.
The same question on the released text scores 72 of 150, and the paired contrast between the two inputs is 52 wins to 17 losses ($p<0.0001$).
Context helps further: $\pm$15 segments against the same span is 30 wins to 10 losses ($p=0.002$).
The word-count rules reach 66.7\% and 68.0\%, but their paired tests against the majority label are not significant ($p=0.10$ and $p=0.053$).

Table~\ref{tab:graph} gives the argument-based methods.
With Terra's arguments and Jev's scores the full factor graph of \citet{wu2025power} reaches 72.0\%.
That is level with Terra asked directly (72.7\%, 14 wins to 15 losses) and with the direct Jev question (14 to 13).
With Haiku's arguments it reaches 70.0\% under Jev's scores and 68.0\% under Haiku's, against 68.0\% for Haiku asked directly.
The parts of the graph do not behave consistently across the two sets of arguments.
Rebuttals add 5.3 points with Terra's arguments (11 wins to 3 losses, $p=0.057$) and cost 3.3 points with Haiku's under the same scorer, and adding the relations changes accuracy in neither Jev-scored run.
A plain sum of reason scores with no graph equals the reasons-only graph in all three rows.
The best single result at this span is a stronger model asked directly (Sonnet, 76.0\%), which is not distinguishable from Jev (27 wins to 20 losses, $p=0.38$) or from Terra (15 to 10, $p=0.42$).

The grounded extension behaves on the rebuilt snippets as it did on the released ones.
Under the proposed relation no reason is ever \textsc{in} in any run, so all 150 decisions are made by the tie-break and the accuracy is that of the sum of reason strengths.
The preference relation gives a real decision: with Jev scores every snippet is decided by \textsc{in} counts, and accuracy is 70.0\% with Terra's arguments and 73.3\% with Haiku's, the latter level with the direct Jev question (13 wins to 10 losses, $p=0.68$).
With the coarser Haiku scores it reaches 64.7\%.
We also searched 288 grounded policies over the Haiku arguments and Jev scores, varying the strength definitions, the rule for when a rebuttal attacks (preference or a threshold on Jev's judgment that it counters its reason), the cross-side rule, and the decision rule.
Choosing the best policy on a development half of the meetings and scoring it on the other half, over 500 random splits, gives 69.7\% on average against 71.5\% for the direct Jev question on the same halves, so no policy in this family improves on asking directly.

The typed-statement graph tracks the direct Jev question at every input (70.0\% against 71.3\% at the same span, 21 wins to 23 losses, $p=0.88$).
It is a cheaper way to run the method and exposes which statements drive a decision, and we do not claim it is more accurate than asking directly.

Jev's probabilities are informative on the rebuilt snippets.
Its Brier score is 0.189 at the same span and 0.095 at $\pm$60, against 0.246 for the base rate, and its most confident fifth of answers is 93.3\% correct at the same span.

\section{Limitations}

The order of work was exploratory.
The mismatch in the released file was found, and the rebuilt set constructed, after many runs on the released file, and every number in Section~\ref{sec:released} was computed on test data with no development split.
The questions used on the rebuilt snippets were written before that data existed, and each run was made once.

Neither rerun of the factor graph uses the original model.
The Terra run follows the pipeline call by call, but the text comes from GPT-6.1 Terra and the scores from Jev, because the command-line interface we used does not expose token probabilities.
The Haiku run wrote each snippet's arguments in one pass, so reasons were written with the conversation in view, and its Haiku scores are stated probabilities with about 15 distinct values.
In both, each rebuttal is judged against its own reason and grounded on its own text.

The rebuilt set has four professors, and one of them is the professor in 110 of the 150 snippets.
Jev's accuracy for that professor is 69.1\% at the same span, and 73.3 to 100\% for the others, so the result is not carried by one person, but the meeting-level bootstrap is the more honest interval.
With 60 extra segments a model can use cues that a six-turn excerpt lacks, including how others address the two speakers.

We did not rerun GPT-3.5, Gemma, or a human study on the rebuilt snippets.
Jev is a proprietary model in early access whose training and calibration claims are not peer reviewed.
The 40-item human audit from our proposal was not completed.

Earlier work predicts power from word choice and thread structure with supervised models \citep{danescu2012echoes,prabhakaran2014power}, and \citet{ziems2024css} find that LLMs often trail such models.
Our numbers say only that the role is recoverable from these excerpts, and they are not a comparison with that line of work.

\section{Reproducibility}

The repository contains the rebuilt snippets with corpus offsets, every model answer, and all code.\footnote{\url{https://github.com/muttaqiahmadalladin/cs-577}}
One command downloads the five external inputs (the released \texttt{meeting.csv}, argument cache and rule file, and the two ICSI annotation archives), verifies each by SHA-256, regenerates the snippet set byte for byte, and recomputes every number in this report from the stored model answers without calling a model.
The model answers themselves come from hosted models and cannot be regenerated exactly; the scripts and instructions used to collect them are included.
The ICSI corpus is distributed under CC BY 4.0.

\bibliography{references}

\end{document}
